\section{Benchmark Design and Reproducibility}
\label{app:benchmarking}

% Question: What detector data form the benchmark, and how is the synthetic transit added?
% Answer: A known transit is projected into a real out-of-transit NIRISS/SOSS time series while preserving the observed background, contaminants, and noise.
% Evidence: Injector manifest and the archived benchmark release.
\subsection{Detector-level benchmark}

The benchmark is constructed from out-of-transit observations of WASP-17 with
NIRISS/SOSS. It contains 229 Stage-2 \texttt{rateints} integrations. Synthetic
transit modulation is applied to integrations 51 through 177, inclusive, so
the injected event contains 127 integrations and the remaining 102 form its
complement. The delivered science values, uncertainties, and data-quality
flags are retained, and no background subtraction is undone or reapplied.

The injection is defined by Equation~\ref{eq:v45-injector}. In that
expression, $Y^{\mathrm{OOT}}_{tp}$ is the observed out-of-transit detector
time series, $\bm{s}_{o,t}$ is the denoised expectation of the target source
for order $o$, $\bm{T}^{\mathrm{true}}_{o,t}$ is the known injected transit
factor, and $\bm{A}_o$ is the corresponding detector projection operator.
Only the target-source expectation is modulated. The realised detector noise,
additive background, field sources, higher-order structure, and temporal
variability remain as observed. The source expectation is denoised before the
transit is applied so that the injector does not attenuate a particular noise
realisation. Contributions from the two spectral orders are constructed
separately and summed in detector space, including where their supports
overlap.

The projection operators use the PASTASOSS trace and wavelength solution
\citep{BainesEtAl2023Trace,BainesEtAl2023Wavelength} and the WebbKernel
implementation of the ATOCA response mapping
\citep{DarveauBernierEtAl2022,JWSTPipeline2025}. If empirical target light
falls outside the valid model response, the injection stops instead of
clipping the source or assigning it a unity transit factor. Unsupported
source flux therefore cannot pass through the injector without modulation.

% Question: How is information about the injected spectrum kept out of the recovery methods?
% Answer: The injected spectrum and generator products are unavailable during recovery, and scoring begins only after the detector prediction is final.
% Evidence: Benchmark manifests and the recovery and scoring interfaces.
\subsection{Blind recovery and scoring}

The injected spectrum, source expectation, and detector projection are not
available to the recovery analyses. Every method instead receives the same
detector cube and declared calibration products, constructs its own response,
masks, nuisance model, and starting points, and produces a final detector
prediction. Only then is the injected spectrum revealed and compared with the
recovery. This separation prevents the truth from influencing calibration,
model selection, convergence, or the reported spectrum.

If inspection of a recovered-versus-injected residual motivates a change to a
method, the revised method is evaluated on a new hidden injection before that
change contributes to an accuracy claim. The same design applies to a future
repair of the known spatial-response defect: a correction family can be
motivated by the observed failure mode, but it will be calibrated on
alternating out-of-transit folds, selected by held-out detector likelihood,
and fixed before comparison with a new hidden spectrum.

% Question: Which wavelengths and metrics enter the common comparison?
% Answer: Every method is scored on complete R=100 bins within the direct PASTASOSS calibration domain.
% Evidence: Calibrated-domain mask and common-grid manifest.
\subsection{Common reporting support}

All methods are compared on the same physical $R=100$ grid and with the same
predefined support mask. The order-1 PASTASOSS training anchors for this visit
end at detector coordinate $x=73.62893$, corresponding to
$2.75977449\,\mu\mathrm{m}$. Detector columns 38 through 73 depend on
polynomial extrapolation and are therefore excluded from the primary
comparison. Any reporting bin with even partial dependence on these columns is
removed in full. This leaves 141 complete bins, with the final retained bin
ending at approximately $2.740018\,\mu\mathrm{m}$. The mask is determined by
the calibration domain, applied identically to all methods, and used only for
reporting rather than fitting.

The primary accuracy statistic is the unweighted absolute root-mean-square
error on these common bins. Mean bias and morphology error are reported
separately so that an achromatic offset is not confused with an incorrect
spectral shape. Order-specific, overlap, crossing-region, and supported-red
diagnostics accompany the aggregate score. Weighted error is treated as a
sensitivity statistic. Zero-injection, flat-depth, and noiseless engineering
tests are reported separately from the main recovery experiment.

% Question: Can every comparison pipeline reproduce its normal preprocessing from an integration-level cube?
% Answer: No. A matched group-stage experiment is required to isolate the effect of preprocessing performed before ramp fitting.
% Evidence: Injector fairness programme and current comparison status.
\subsection{Integration-level limitation and paired group-stage experiment}

The present benchmark begins with integration-level count-rate images. It does
not contain the nondestructive detector groups from which those images were
fitted. Pipelines such as Ahsoka and supreme-SPOON/exoTEDRF normally perform
characteristic $1/f$ corrections before or during ramp fitting and cannot
reproduce that step from the delivered cube. The current benchmark can compare
recoveries from identical integration-level data, but it cannot determine how
group-stage preprocessing changes their relative performance.

A paired sensitivity experiment is therefore planned. One common group-level
baseline will be used to create matched group-stage and integration-stage
injections from the same physical source deficit. Each pipeline will retain
its own predefined preprocessing, and NOVA will also be run after the same
group-stage $1/f$ correction. The comparison will include zero-injection
identity, the delivered detector deficit, white-light scatter, red-end
recovery, and common-bin spectral error. This experiment does not form part of
the present result.

A preliminary scene-aware injector is also excluded from the current
benchmark because its broad full-column masks remove a material fraction of
the target signal together with contaminants. A future version will require
out-of-transit scene modelling that protects the target trace and quantifies
any remaining contaminant leakage.

% Question: What additional experiments are needed before the comparison can be generalised?
% Answer: Performance must remain stable across detector-response mismatch, visits, atmospheric spectra, and noise realisations.
% Evidence: Benchmark limitations and planned validation programme.
\subsection{Structural neutrality and ensemble tests}

The injector and NOVA use related families of trace geometry, wavelength
mapping, and spectral mixing. They do not share fitted detector responses or
the same executable implementation, but their common modelling ancestry
creates an inverse-crime risk. The benchmark programme will therefore include
hidden injections with deliberately mismatched spatial profiles, wavelength
maps, line-spread functions, throughput, background morphology, and temporal
systematics.

WASP-17 is a development visit rather than sufficient validation by itself.
Additional visits will be assessed for target and visit identity, observing
mode, out-of-transit coverage, integration chronology, data quality,
background structure, field sources, and order-3 morphology. The same injected
spectrum will first be recovered across suitable visits, followed by
additional spectral morphologies and a method-neutral noise ensemble. At least
one visit will remain outside model-branch selection. Until these tests and
the comparison-pipeline matrix are complete, the present benchmark supports a
controlled single-visit result rather than a universal ranking of reduction
methods.

% Question: What information accompanies a reported benchmark recovery?
% Answer: The inputs, fitted configuration, convergence evidence, spectrum, covariance, scores, and detector diagnostics needed to reproduce and interpret the result.
% Evidence: Release manifest and portable handoff audit.
\subsection{Reproducibility record}

For each recovery, the archived record identifies the detector cube,
uncertainties, flags, time grid, wavelength support, code, configuration,
calibration and response products, and final detector prediction. It also
contains the termination state, final-weight confirmation, multistart
agreement, numerical conditioning, active bounds, recovered spectrum and
covariance, common-grid accuracy scores, and residual and robust-weight
diagnostics by time, order, and detector region. Content hashes link large
cluster-resident arrays to the portable record. Together these products make
the reported recovery reproducible and allow its spectral accuracy to be
interpreted alongside the detector fit that produced it.
